\begin{tikzpicture}[
    font=\footnotesize,
    >=Stealth,
    % ---------------------------------------------------------------
    % Global Styling & Dimensions
    % ---------------------------------------------------------------
    baseNode/.style={
        rectangle,
        draw=black!80,
        thick,
        rounded corners=4pt,
        align=center,
        inner sep=7pt
    },
    wideNode/.style={
        baseNode,
        text width=12.2cm
    },
    stageNode/.style={
        baseNode,
        fill=blue!5,
        text width=12.2cm,
        minimum height=1.15cm
    },
    evalBranch/.style={
        baseNode,
        fill=blue!10,
        text width=5.6cm,
        minimum height=1.15cm
    },
    arrow/.style={
        ->,
        >=Stealth,
        thick,
        draw=black!75
    },
    branchArrow/.style={
        ->,
        >=Stealth,
        thick,
        draw=black!65
    }
]

    % ==================================================================
    % 1. EXPERIMENTAL INPUT
    % ==================================================================
    \node[wideNode, fill=yellow!15] (problem) {
        \textbf{Optimization Problem and Experimental Setup} \\[2pt]
        \scriptsize
        Continuous black-box optimization \quad $\bullet$ \quad
        BBOB benchmark functions \quad $\bullet$ \quad
        Dimension $D$ \quad $\bullet$ \quad
        Evaluation budget $B$
    };

    % ==================================================================
    % 2. UNIFIED PROMPT SCAFFOLDING BLOCK
    % ==================================================================
    \node[
        wideNode,
        fill=orange!10,
        draw=orange!80!black,
        below=0.9cm of problem
    ] (promptScaffolding) {
        \textbf{Prompt Scaffolding Framework} \\[2pt]
        \scriptsize
        Invariant layout, API templates, and defensive constraints \\
        coupled with $3 \times 4$ factorial synthesis conditions:
        information modes ($\textsc{Clean}, \textsc{Implicit}, \textsc{Noisy}$)
        $\times$ search strategy scaffolds
    };

    % ==================================================================
    % 3. LLAMEA SYNTHESIS
    % ==================================================================
    \node[
        stageNode,
        below=0.9cm of promptScaffolding
    ] (synthesis) {
        \textbf{LLaMEA Algorithm Synthesis} \\[2pt]
        \scriptsize
        Iterative LLM-driven generation, validation, execution, and semantic selection of candidate heuristics
    };

    % ==================================================================
    % 4. SYNTHESIS OUTPUT
    % ==================================================================
    \node[
        stageNode,
        below=0.85cm of synthesis
    ] (champions) {
        \textbf{Synthesized Algorithms} \\[2pt]
        \scriptsize
        Retained candidate optimization algorithms and corresponding generational telemetry
    };

    % ==================================================================
    % 5. INDEPENDENT EVALUATION HARNESS
    % ==================================================================
    \node[
        stageNode,
        below=0.85cm of champions
    ] (evaluation) {
        \textbf{Independent Algorithm Evaluation} \\[2pt]
        \scriptsize
        Systematic validation of synthesized heuristics across target benchmark environments
    };

    % ==================================================================
    % 6. DUAL EVALUATION BRANCHES (CLEAN & NOISY)
    % ==================================================================
    \node[
        evalBranch,
        below left=0.8cm and 0.25cm of evaluation.south,
        anchor=north east
    ] (cleanEvaluation) {
        \textbf{Clean Evaluation} \\[2pt]
        \scriptsize
        Performance on underlying deterministic BBOB objectives
    };

    \node[
        evalBranch,
        below right=0.8cm and 0.25cm of evaluation.south,
        anchor=north west
    ] (noisyEvaluation) {
        \textbf{Noisy Evaluation} \\[2pt]
        \scriptsize
        Performance under heteroscedastic stochastic observations
    };

    % ==================================================================
    % 7. STATISTICAL ANALYSIS
    % ==================================================================
    \coordinate (evalMid) at ($(cleanEvaluation.south)!0.5!(noisyEvaluation.south)$);

    \node[
        wideNode,
        fill=purple!10,
        below=0.85cm of evalMid
    ] (analysis) {
        \textbf{Statistical Analysis and Comparison} \\[2pt]
        \scriptsize
        Target precision trajectories, empirical runtime distributions (ECDF), and reference benchmark comparisons
    };

    % ==================================================================
    % 8. FINAL RESULTS
    % ==================================================================
    \node[
        wideNode,
        fill=green!15,
        below=0.85cm of analysis
    ] (results) {
        \textbf{Experimental Results and Research-Question Analysis} \\[2pt]
        \scriptsize
        Comparative assessment against CMA-ES, DE, and PSO to address core research questions
    };

    % ==================================================================
    % PIPELINE CONNECTORS
    % ==================================================================
    \draw[arrow] (problem) -- (promptScaffolding);
    \draw[arrow] (promptScaffolding) -- (synthesis);
    \draw[arrow] (synthesis) -- node[fill=white, inner sep=2pt, font=\scriptsize\bfseries] {Synthesized Candidates} (champions);
    \draw[arrow] (champions) -- (evaluation);

    % Dual Evaluation Forks
    \draw[branchArrow] (evaluation.south) -- ++(0,-0.3) -| (cleanEvaluation.north);
    \draw[branchArrow] (evaluation.south) -- ++(0,-0.3) -| (noisyEvaluation.north);

    % Dual Evaluation Convergence into Analysis
    \draw[branchArrow] (cleanEvaluation.south) -- ++(0,-0.3) -| ([xshift=-3.0cm]analysis.north);
    \draw[branchArrow] (noisyEvaluation.south) -- ++(0,-0.3) -| ([xshift=3.0cm]analysis.north);

    \draw[arrow] (analysis) -- (results);

    % ==================================================================
    % STAGE GROUPING CONTAINERS (BACKGROUND)
    % ==================================================================
    \begin{scope}[on background layer]
        % Synthesis Stage Container
        \node[
            draw=blue!50!black,
            dashed,
            fill=blue!2,
            rounded corners=6pt,
            inner sep=9pt,
            fit=(synthesis) (champions),
            label={[font=\scriptsize\bfseries, text=blue!60!black]left:Synthesis}
        ] (synthesisGroup) {};

        % Independent Evaluation Container
        \node[
            draw=blue!50!black,
            dashed,
            fill=blue!2,
            rounded corners=6pt,
            inner sep=9pt,
            fit=(evaluation) (cleanEvaluation) (noisyEvaluation),
            label={[font=\scriptsize\bfseries, text=blue!60!black]left:Evaluation}
        ] (evaluationGroup) {};

        % Analysis Stage Container
        \node[
            draw=purple!60!black,
            dashed,
            fill=purple!2,
            rounded corners=6pt,
            inner sep=9pt,
            fit=(analysis) (results),
            label={[font=\scriptsize\bfseries, text=purple!70!black]left:Analysis}
        ] (analysisGroup) {};
    \end{scope}

\end{tikzpicture}